% !TeX root = ../nc-quadric-classification.tex
\section{階数三支持超平面の共通二変数環}
\label{sec:common-plane}

基礎体 $k$ は標数零の代数閉体とする。
次数一で生成された連結二次次数付き domain $B=T(V)/(K)$ について
\[
 \dim V=4,\qquad K\subset V\otimes V,\qquad \dim K=7
\]
を仮定する。従って $\dim B_2=9$ である。
線形行列 $M\in\operatorname{Mat}_2(V)$ の成分が張る部分空間を $U(M)$ と書く。

\begin{lemma}[行列積の四成分の独立性]
\label{lem:w-matrix-products}
$L_1,L_2\subset\operatorname{Mat}_2(k)$ を三次元部分空間とする。
trace pairing に関する $L_2^\perp$ が可逆行列で生成されるなら、
行列の乗法
\[
 L_1\otimes L_2\longrightarrow\operatorname{Mat}_2(k)
\]
は全射である。
従って $P\in\operatorname{Mat}_2(U_1)$ と $Q\in\operatorname{Mat}_2(U_2)$ の係数行列が
それぞれ $L_1,L_2$ の基底をなすとき、$PQ$ の四成分は
形式的テンソル空間 $U_1\otimes U_2$ で一次独立である。
\end{lemma}
\begin{proof}
全射でないなら、非零行列 $D$ があって
$\operatorname{tr}(DXY)=0$ が全ての $X\in L_1,Y\in L_2$ に対して成立する。
従って $DL_1\subset L_2^\perp=kC$ であり、$C$ は可逆である。
$D$ が可逆なら $\dim DL_1=3$ なので矛盾する。
$\operatorname{rank} D=1$ なら $DX$ は全て特異行列であるため $DL_1=0$ となるが、
左乗法 $X\mapsto DX$ の核は二次元であり、三次元の $L_1$ を含めない。
後半はこの全射の双対を取ったものである。
\end{proof}

\begin{lemma}[平方零行列の成分の関係]
\label{lem:w-invertible-annihilator}
$M\in\operatorname{Mat}_2(V)$ が $B$ 上で $M^2=0$ を満たし、$\dim U(M)=3$ とする。
四成分の唯一の一次関係を表す二次正方係数行列は可逆である。
従って $M^2$ の四成分を $V\otimes V$ の元として解釈すると、
それらは四次元部分空間 $R(M)\subset K\cap U(M)^{\otimes2}$ を張る。
\end{lemma}
\begin{proof}
関係の係数行列の階数が一なら、非零の定数行ベクトル $a$ と
列ベクトル $b$ によって $aMb=0$ と書ける。
$ab\ne0$ なら同時基底変換、すなわち $M$ の共役変換によって
この式を $M_{11}=0$ にできる。
このとき $M=\left(\begin{smallmatrix}0&x\\y&z\end{smallmatrix}\right)$
と書け、$M^2=0$ の第一対角成分から $xy=0$ となる。
domain 性により $x=0$ または $y=0$ で、成分空間の次元三に反する。
$ab=0$ なら同時基底変換で $M_{12}=0$ にできる。
このとき $M=\left(\begin{smallmatrix}x&0\\y&z\end{smallmatrix}\right)$
であり、二つの対角成分から $x^2=z^2=0$ となって同様に矛盾する。
関係の係数行列の転置は係数行列空間の trace-annihilator を生成するので、
補題~\ref{lem:w-matrix-products} を適用できる。
\end{proof}

\begin{lemma}[二次元部分空間にある関係]
\label{lem:w-binary-relations}
二次元部分空間 $W\subset V$ に対し、$\dim(K\cap W^{\otimes2})\le1$ である。
等号が成り立つ場合、非零の関係は二因子テンソルとして階数二である。
\end{lemma}
\begin{proof}
階数一の非零テンソル $w\otimes w'$ が $K$ にあれば、
$B$ で $ww'=0$ となって domain 性に反する。
もし $K\cap W^{\otimes2}$ に独立な二つの元 $r_0,r_1$ があれば、
$\det(sr_0+tr_1)$ は $(s,t)$ の斉次二次式である。
代数閉性により、この式が恒等的に零の場合も含めて、
非零の $(s,t)$ で行列式が零になる。
$sr_0+tr_1$ は非零なので階数一となり、矛盾する。
\end{proof}

\begin{theorem}[non-standard 二型の共通二変数環]
\label{thm:w-common-plane}
$M_1,M_2\in\operatorname{Mat}_2(V)$ が $B$ 上で $M_i^2=0$ を満たすとする。
$U_i=U(M_i)$ は三次元で、$U_1\ne U_2$ と仮定する。
$W=U_1\cap U_2$ および $R_i=R(M_i)$ と置けば、
\[
 \dim W=2,\qquad
 R_1\cap R_2=K\cap W^{\otimes2}=kr,\qquad
 \operatorname{rank} r=2,\qquad K=R_1+R_2.
\]
特に二つの平方零行列は $B$ の全ての二次関係を与える。
\end{theorem}
\begin{proof}
補題~\ref{lem:w-invertible-annihilator} により $\dim R_i=4$ である。
$R_i\subset K$ と $\dim K=7$ から $\dim(R_1\cap R_2)\ge1$ となる。
$U_1,U_2$ に適合する $V$ の基底を取れば
\[
 (U_1\otimes U_1)\cap(U_2\otimes U_2)
   =(U_1\cap U_2)\otimes(U_1\cap U_2)=W^{\otimes2}
\]
である。従って $R_1\cap R_2\subset K\cap W^{\otimes2}$ であり、
補題~\ref{lem:w-binary-relations} によって両辺は一次元に等しい。
さらに $\dim(R_1+R_2)=4+4-1=7=\dim K$ となる。
\end{proof}

\paragraph{正規切断への適用.}
系~\ref{cor:nonstandard-square-matrix} により、
正規切断 $B=A/(g)$ で $p(X)=X,p(Y)=Y$ のとき、評価行列は
$M_X^2=gC_X$, $M_Y^2=gC_Y$ ($C_X,C_Y\in\operatorname{GL}_2(k)$) と取れる。
$B$ に移せば平方零なので、双方が階数三かつ支持が異なる場合には
定理~\ref{thm:w-common-plane} が適用される。
$g$ が中央であることも $C_X,C_Y$ がスカラーであることも使用していない。

\begin{remark}[交差積の場合にも使える部分]
\label{rem:w-crossed-products}
二つの行列 $P,Q$ の支持が異なる三次元空間 $U_1,U_2$ で、
それぞれの係数行列空間の trace-annihilator が可逆行列で生成され、
$PQ=QP=0$ が $B$ で成立する場合も同じ結論が成り立つ。
この場合、四次元の関係空間は
$R_{12}\subset U_1\otimes U_2$ と $R_{21}\subset U_2\otimes U_1$ にあり、
それらの交わりが $W^{\otimes2}$ に含まれることを使う。
中央切断の通常の行列分解 $PQ=QP=gI$ はこの状況の候補となるが、
一般の正規切断では一方の積に正規化自己同型が入る。
後者への適用には、twist された支持の一致を別に確定する必要がある。
\end{remark}

\subsection{全次数での量子平面・Jordan 平面の埋込み}

\begin{proposition}
\label{prop:w-embedded-plane}
$W\subset B_1$ が二次元で $K\cap W^{\otimes2}=kr\ne0$ とする。
このとき $W$ の基底 $u,v$ を適切に選ぶと、$W$ が生成する $B$ の部分環は
次のいずれかと同型である。
\[
 Q_q=k\langle u,v\rangle/(vu-quv),\quad q\in k^\times,
 \qquad
 Q_J=k\langle u,v\rangle/(vu-uv-u^2).
\]
これは二次関係だけの同定でなく、自然な写像 $Q_q\to B$ または $Q_J\to B$
の全次数での単射を主張する。
\end{proposition}
\begin{proof}
補題~\ref{lem:w-binary-relations} により $r$ の係数行列は可逆である。
$r$ を対称部分と反対称部分に分ける。
対称部分の階数が二なら、代数閉体上の二変数非退化対称形式を
$uv+vu$ に正規化できる。従って $r=\alpha uv+\beta vu$ と書け、
係数行列の可逆性から $\alpha\beta\ne0$ である。これは $Q_q$ の関係である。
対称部分が零なら $r$ は $vu-uv$ の定数倍である。
対称部分の階数が一なら、$r=\alpha u^2+\beta(vu-uv)$ と書ける。
可逆性から $\beta\ne0$ なので、基底を定数倍して Jordan 平面の関係となる。

次に単射性を証明する。いずれの環 $Q$ も $k[u]$ の Ore 拡大で、
単項式 $u^iv^j$ ($i,j\ge0$) が基底をなし、domain である。
$u$ は正規なので可逆化できる。
$z=u^{-1}v\in Q[u^{-1}]$ と置けば、
\[
 zu=u\tau(z),\qquad
 \tau(z)=qz\ \ (Q=Q_q),\qquad
 \tau(z)=z+1\ \ (Q=Q_J)
\]
となる。従って次数 $n$ の任意の非零斉次元 $f$ は一意に
\[
 f=u^n p(z),\qquad 0\ne p\in k[z],\qquad \deg p\le n
\]
と書ける。実際、$u^{n-j}v^j$ をこの形に直すと、
$p$ は次数 $j$ の非零先頭係数を持つ多項式になり、三角的に基底を与える。
代数閉性から $p(z)=a\prod_{i=1}^d(z-\lambda_i)$ と書く。
恒等式
\[
 \prod_{i=1}^{d}\bigl(u f_i(z)\bigr)
 =u^d\prod_{i=1}^{d}\tau^{d-i}\bigl(f_i(z)\bigr)
\]
に $f_i(z)=\tau^{-(d-i)}(z-\lambda_i)$ を入れれば、
$u^dp(z)$ は $Q_1$ の非零元 $d$ 個の積の定数倍となる。
各 $f_i$ が一次式なので $u f_i(z)\in ku+kv$ である。
残る $u^{n-d}$ も合わせると、$f$ は非零一次元 $n$ 個の積である。

自然な全射 $Q\to k\langle W\rangle\subset B$ は次数一で単射である。
その核が非零なら非零斉次元 $f$ を取り、上の一次因子分解を写すと、
$B$ で非零一次元の積が零になる。これは domain 性に反する。
\end{proof}

\begin{remark}[ここからはまだ正規化できない]
$Q_q$ の $u,v$ は $Q_q$ の中で正規であり、$Q_J$ の $u$ は $Q_J$ の中で正規である。
しかし $B$ 内での正規性は、この埋込みの命題からは従わない。
また $q\ne1$ の量子平面や Jordan 平面の場合、$W$ の線形基底変更だけで
二つの生成元を $B$ の中心に入れることはできない。
そのような基底があれば $W$ の全ての元が互いに交換するが、
命題の全次数単射に反するからである。
Zhang twist や pointed 同値によって別の表示へ移る可能性は残る。
\end{remark}

\subsection{同一支持と中央化についての障害}

\begin{lemma}[同一支持の場合に確定すること]
\label{lem:w-equal-support}
三次元部分空間 $U\subset B_1$ について
$\dim(K\cap U^{\otimes2})\le4$ である。
従って定理~\ref{thm:w-common-plane} の平方零行列で $U_1=U_2=U$ なら、
\[
 R_1=R_2=K\cap U^{\otimes2}.
\]
\end{lemma}
\begin{proof}
$\mathbb P(U)\times\mathbb P(U)$ の Segre 像は
$\mathbb P(U^{\otimes2})=\mathbb P^8$ 内の四次元射影多様体である。
$K\cap U^{\otimes2}$ が五次元以上なら、その射影化の次元は四以上であり、
射影次元定理によりこの Segre 像と交わる。
この交点は非零の階数一テンソルを与え、domain 性に反する。
後半は各 $R_i$ が四次元であることから従う。
\end{proof}

\paragraph{残る点.}
$R_1=R_2$ から二つの余核が同型になるかは、この次元論だけでは分からない。
従って二つの異なる MCM 型で支持が異なるという仮定を外してはいけない。
次の具体的課題は、三次元の同じ成分空間と同じ四次元の平方関係空間を持つ
二つの行列について、定数行列による左右同値を証明するか、反例を得ることである。

\paragraph{中央化は適合表示の中ですら自動でない.}
命題~\ref{prop:nonstandard-rank-three-example} で使う
\[
 B_0=k[u,v]\langle c,d\rangle/(c^2-uv,d^2-uv)
\]
を考える。この表示では $u,v$ は中央である。
$t\in k^\times$ に対して
\[
 \sigma_t(u)=t^2u,\quad \sigma_t(v)=v,\quad
 \sigma_t(c)=tc,\quad\sigma_t(d)=td
\]
は次数付き自己同型である。
$a\star b=a\sigma_t^{\deg a}(b)$ の規約の Zhang twist を $B_t$ と書くと、
関係は
\[
\begin{gathered}
 vu=t^2uv,\qquad cu=tuc,\qquad cv=t^{-1}vc,\\
 du=tud,\qquad dv=t^{-1}vd,\qquad
 c^2=tuv,\qquad d^2=tuv.
\end{gathered}
\]
ここでは積記号を再び通常の並置で書いた。
行列
\[
 M_r=\begin{pmatrix}r&u\\-tv&-tr\end{pmatrix},\qquad r=c,d,
\]
はそれぞれ $M_r^2=0$ を満たす。
この二つの支持は $\langle u,v,c\rangle$, $\langle u,v,d\rangle$ であり、
交わりはパラメータ $q=t^2$ の量子平面を生成する。
$t^2\ne1$ なら、同じ $B_t$ の中で $W$ の基底を変えて中央化することは不可能である。

特に $t=\mathrm i$ は命題~\ref{prop:nonstandard-rank-three-example} の反 Serre 平方根に適合した表示そのものであり、
このとき $vu=-uv$ である。
従って「$H_2$ の適合表示なら $W$ を線形基底変更だけで中央化できる」
という主張は、既存例によって既に否定される。
一方 $u,v$ はこの例では $B_t$ 全体でも正規なので、
一般の $W$ の正規化という課題の反例にはなっていない。
任意の $t$ について次数シフト自体が反 Serre の平方根になるとは主張していない。
